\documentclass[../DoAn.tex]{subfiles}
\begin{document}

\begin{center}
    \Large{\textbf{ABSTRACT}}\\
\end{center}
\vspace{1cm}

Quantized artificial intelligence is an important approach for deploying inference models on resource-constrained embedded systems. On small RISC-V processors, common INT8 operations such as multiply-accumulate, dot product, activation, and value clipping are often expressed as multiple scalar instructions, increasing the dynamic instruction count and loop-control overhead. Software optimization, SIMD/DSP support, RISC-V Vector (RVV), and dedicated accelerators can all improve performance, but they also introduce trade-offs in hardware cost, flexibility, and system integration. Therefore, this thesis adopts an instruction-set-level approach, aiming to reduce instruction count and control overhead directly within the processor execution flow instead of offloading the entire computation to a separate accelerator.

Based on this approach, the thesis specifies a set of operations for quantized AI kernels, including multiply-accumulate, packed dot product, activation, value clipping, post-increment load, and hardware loop. These operations are modeled in gem5 to validate their semantics and study performance trends. In parallel with simulation, a minimal SoC based on the CV32E40P core is implemented on the DE2i-150 FPGA board. The corresponding instruction extension is configured on the core, and firmware is developed to create evaluation scenarios, including an INT8 MNIST fully connected model using TensorFlow Lite Micro. The results show that optimized kernels in simulation can achieve over \(9\times\) speedup under a suitable configuration, while the FPGA implementation reaches about \(12\times\)--\(13\times\) speedup over the software reference in the main experiments. The main contribution of this thesis is not only the design of an extended RISC-V instruction set for quantized AI inference, but also the establishment of a continuous validation flow from simulation to real hardware implementation. This demonstrates the potential of the proposed approach on resource-constrained embedded platforms.

\end{document}
